\begin{proof}[Proof of \cref{obj:lem:cty-winsorized-score-maximal-bound}]
Throughout, \(Z=(Y(0),Y(1),X)\) denotes one potential-outcome observation, \(Y=TY(1)+(1-T)Y(0)\) its observed outcome, and for \(u\in\mathbb R^{p+1}\) we write \(\lVert u\rVert=\max_{0\leq j\leq p}|u_j|\) for the coordinate maximum, which is the norm appearing in the displayed conclusion.  All constants below are fixed before the moment exponent \(\nu\), the law \(P\), the sample size \(n\), the bandwidth \(h\), and the winsorization level \(B\).

\textbf{Step 1: a coefficient radius depending only on \(p\) and \(L\).}  Set
\[
R=1+\left|L(p+1)\cdot 16L(2+L)\right|,
\]
which is positive.  We claim that for every \(P\in\mathcal P_{12}(p,\nu,L)\), every \(t\in\{0,1\}\), every \(x\in\mathcal B_P\), every \(0<h\leq L^{-1}\) and every \(0\leq j\leq p\),
\[
\left|\beta^*_{t,P,x}(h)_j\right|\leq R .
\]
Write
\[
m_{t,P,x}(h)
=
\mathbb E_P\!\left[
{1\over h^2}\mathbf 1\{D^{\pm}_{P,x}\in I_t\}
K_\square(D^{\pm}_{P,x}/h)\,r_p(D^{\pm}_{P,x}/h)\,Y
\right]
\]
for the right-hand side of the population normal equations in \cref{obj:synth_128}.  Its coordinates obey
\[
\max_{0\leq j\leq p}\left|m_{t,P,x}(h)_j\right|\leq 16L(2+L).
\]
Indeed, the uniform kernel forces \(|D^{\pm}_{P,x}|\leq h\), hence \(X\in B_2(x,h)\) and \(|r_p(D^{\pm}_{P,x}/h)_j|\leq1\), so the integrand is dominated pointwise by \(h^{-2}\mathbf 1\{X\in B_2(x,h)\}\{1+2Y(0)^2+2Y(1)^2\}\).  The density envelope \(f_P\leq L\) of \cref{obj:def:cty-a1-a2-class} and the area \(\pi h^2\leq4h^2\) of a planar disk give
\[
P\{X\in B_2(x,h)\}\leq 4Lh^2 .
\]
The moment envelope in \cref{obj:def:cty-a1-a2-class} gives, since \(\nu\geq2\), the conditional second-moment estimate
\[
\mathbb E_P\bigl[Y(t)^2\mid X=x\bigr]\leq 1+L,
\qquad x\in\mathcal X_P,
\]
using \(y^2\leq 1+|y|^{2+\nu}\).  Localizing this bound to \(B_2(x,h)\) gives \(\mathbb E_P[\mathbf 1\{X\in B_2(x,h)\}Y(t)^2]\leq (1+L)4Lh^2\) for each arm, so the displayed coordinate bound follows from
\[
h^{-2}\bigl\{4Lh^2+4(1+L)4Lh^2\bigr\}
=20L+16L^2
\leq
16L(2+L).
\]
% lean: populationScore_apply_uniform_bound
% lean: marginal_closedBall_le
% lean: localized_arm_sq_lintegral_le
The population Gram lower bound of \cref{obj:def:cty-a1-a2-class} and the normal equations of \cref{obj:synth_128}, written with \(\beta^*=\beta^*_{t,P,x}(h)\), then give
\[
L^{-1}\lVert\beta^*\rVert^2
\leq
L^{-1}\sum_{k=0}^p(\beta^*_k)^2
\leq
(\beta^*)^\top\Psi_{t,P,x}(h)\beta^*
=
(\beta^*)^\top m_{t,P,x}(h)
\leq
(p+1)\lVert\beta^*\rVert\cdot 16L(2+L),
\]
the first step because \(\lVert\beta^*\rVert^2\) is one of the summands.  Hence \(\lVert\beta^*\rVert\leq L(p+1)\cdot16L(2+L)\leq R\), which is the claim.
% lean: populationCoefficient_uniform_bound_explicit

\textbf{Step 2: the bounded score class.}  For \(b\in\mathbb R^{p+1}\) let \(\Pi_R b\) be its componentwise clipping,
\[
(\Pi_R b)_k=\max\{-R,\min\{R,b_k\}\},\qquad 0\leq k\leq p,
\]
and let
\[
\mathcal I_{P,h}
=
\left\{
\iota=(q,t,x,b,j):
q\in[h,2h),\ t\in\{0,1\},\ x\in\mathcal B_P,\ b\in\mathbb R^{p+1},\ 0\leq j\leq p
\right\}.
\]
For \(\iota=(q,t,x,b,j)\in\mathcal I_{P,h}\) define the scalar function of one observation
\[
g_\iota(Z)
=
\mathbf 1\{D^{\pm}_{P,x}\in I_t\}
K_\square\!\left({D^{\pm}_{P,x}\over q}\right)
r_p\!\left({D^{\pm}_{P,x}\over h}\right)_j
\left\{
\psi_B(Y)-r_p\!\left({D^{\pm}_{P,x}\over h}\right)^\top\Pi_R b
\right\}.
\]
Only the kernel window uses the enlarged bandwidth \(q\); the polynomial arguments keep \(h\), so the target score is the member with \(q=h\).  The half-open enlargement \([h,2h)\) supplies the support slack used in Step 6, and clipping the coefficient gives a constant-envelope class.
% lean: separableWinsorizedScoreFunction

\textbf{Step 3: the score is a member of the class.}  Fix \(t\in\{0,1\}\), \(x\in\mathcal B_P\) and \(0\leq j\leq p\), and put \(\iota=(h,t,x,\beta^*_{t,P,x}(h),j)\in\mathcal I_{P,h}\).  By Step 1 the clipping is inactive at this index, \(\Pi_R\beta^*_{t,P,x}(h)=\beta^*_{t,P,x}(h)\), so comparing with the definition of the winsorized score in \cref{obj:synth_131},
\[
\operatorname{score}^{B}_{t,x,h,j}
=
{1\over h^{2}}\cdot{1\over n}\sum_{i=1}^n g_\iota(Z_i).
\]
% lean: winsorizedCenteredScore_apply_eq_separableAverage
The envelope bound of Step 5 makes \(g_\iota\) integrable, so averaging the previous display over the independent and identically distributed observations gives
\[
\mathbb E_{P^n}\!\left[\operatorname{score}^{B}_{t,x,h,j}\right]
=
{1\over h^{2}}\int g_\iota\,dP .
\]
% lean: winsorizedCenteredScore_integral_eq

\textbf{Step 4: reduction to one scalar empirical process.}  For each fixed sample the coordinate maximum \(\lVert\cdot\rVert\) is attained at some \(j\), and every pair \((t,x)\) with \(x\in\mathcal B_P\) together with that \(j\) produces an index of \(\mathcal I_{P,h}\) by Step 3.  Hence, pointwise on the sample space,
\[
\sup_{t\in\{0,1\},\,x\in\mathcal B_P}
\left\lVert
\operatorname{score}^{B}_{t,x,h}
-\mathbb E_{P^n}\!\left[\operatorname{score}^{B}_{t,x,h}\right]
\right\rVert
\leq
{1\over h^{2}}
\sup_{\iota\in\mathcal I_{P,h}}
\left|
{1\over n}\sum_{i=1}^n g_\iota(Z_i)-\int g_\iota\,dP
\right| .
\]
% lean: winsorizedScoreDeviation_le_empiricalProcessSup
Outer expectation is monotone and pulls out the positive finite constant \(h^{-2}\), so
\[
\mathbb E_{P^n}^{*}\!\left[
\sup_{t\in\{0,1\},\,x\in\mathcal B_P}
\left\lVert
\operatorname{score}^{B}_{t,x,h}
-\mathbb E_{P^n}\!\left[\operatorname{score}^{B}_{t,x,h}\right]
\right\rVert
\right]
\leq
{1\over h^{2}}\,
\mathbb E_{P^n}^{*}\!\left[
\sup_{\iota\in\mathcal I_{P,h}}
\left|
{1\over n}\sum_{i=1}^n g_\iota(Z_i)-\int g_\iota\,dP
\right|
\right].
\]
% lean: outerLIntegral_mono
% lean: outerLIntegral_const_mul_le

\textbf{Step 5: envelope, radius and covering witnesses.}  Put
\[
\Lambda_{\mathrm{env}}=(2\cdot4^p+2)\{1+(p+1)R\},
\qquad
\Lambda_{\mathrm{rad}}=1+\left|3\cdot4^p\left\{2(1+L)+\{(p+1)2^pR\}^2\right\}16L\right|,
\]
\[
C_0=1+\max\{\Lambda_{\mathrm{rad}},\Lambda_{\mathrm{env}}\}>0 .
\]
There are constants \(A\geq\exp(1)\) and \(v\geq1\), depending only on \(p\) and \(R\) and therefore only on \(p\) and \(L\), such that for every real \(\nu\), every \(P\in\mathcal P_{12}(p,\nu,L)\) and all \(0<h<B\) with \(B\geq1\), the class \(\{g_\iota:\iota\in\mathcal I_{P,h}\}\) satisfies the following three properties.
\begin{itemize}
\item \textbf{(Envelope.)} For every \(\iota\in\mathcal I_{P,h}\) and every observation \(z\),
\[
|g_\iota(z)|\leq C_0B .
\]
On the kernel window \(|D^{\pm}_{P,x}|\leq q<2h\), so \(|r_p(D^{\pm}_{P,x}/h)_k|\leq2^p\) for every \(k\).  Together with \(|\psi_B(Y)|\leq B\) and \(|(\Pi_Rb)_k|\leq R\), this gives
\[
|g_\iota(z)|
\leq
(2\cdot4^p+2)\{B+(p+1)R\}
\leq \Lambda_{\mathrm{env}}B
\leq C_0B,
\]
where the middle inequality uses \(B\geq1\).
\item \textbf{(Population radius.)} For every \(\iota\in\mathcal I_{P,h}\),
\[
\left(\int g_\iota^2\,dP\right)^{1/2}\leq C_0h .
\]
On the almost-sure support of \(X\), the squared separable score is bounded pointwise by
\[
3\cdot4^p\,\mathbf 1\{X\in B_2(x,2h)\}
\left\{Y(0)^2+Y(1)^2+\{(p+1)2^pR\}^2\right\}.
\]
The kernel window gives the radius \(2h\) because \(q<2h\).  The probability of this disk is at most \(16Lh^2\), while the two localized arm-square integrals are each bounded by \((1+L)16Lh^2\), using the conditional second-moment estimate from Step 1.  Hence \(\int g_\iota^2\,dP\) is bounded by
\[
3\cdot4^p\left\{2(1+L)+\{(p+1)2^pR\}^2\right\}16Lh^2.
\]
Taking square roots and using \(\sqrt u\leq 1+u\) for \(u\geq0\) gives the displayed \(\Lambda_{\mathrm{rad}}h\leq C_0h\) bound.  The winsorization level is absorbed through \(|\psi_B(Y)|\leq |Y|\), so this radius is uniform in \(B\).
\item \textbf{(Polynomial covering.)} For every non-empty finite list \(z_1,\ldots,z_m\) of observations and every \(0<\varepsilon\leq1\) there is a finite set \(\mathcal C\subseteq\mathcal I_{P,h}\) with
\[
\min_{\kappa\in\mathcal C}
\left({1\over m}\sum_{i=1}^m\{g_\iota(z_i)-g_\kappa(z_i)\}^2\right)^{1/2}
<\varepsilon\,C_0B
\quad\text{for every }\iota\in\mathcal I_{P,h},
\]
and with \(|\mathcal C|\leq(A/\varepsilon)^v\).  The certificate is obtained as follows.  First, the fixed-dimensional radial residual-polynomial class with bounded response \(|\psi_B(Y)|\leq B\), bounded arm weights, and coefficient box \([-R,R]^{p+1}\) has polynomial empirical \(L^2\) covers with envelope \(B+(p+1)R\); this is applied with radial ranges \([0,2]\) for the positive and negative arm pieces and \([0,0]\) for the zero-distance and outside-support pieces.  Second, \cref{obj:lem:euclidean-balls-vc} bounds the shattering of the planar closed-ball class, so the classifier family \(\mathbf 1\{\lVert X-x\rVert_2\leq q\}\) indexed by the center-radius parameter has Vapnik--Chervonenkis dimension at most three; the standard passage from a finite Vapnik--Chervonenkis dimension to empirical \(L^2\) covering numbers \citep{Pollard1984,VanderVaart1996} then supplies a fixed polynomial cover for that classifier family.  Product closure combines each radial piece with the ball classifier for the positive and negative pieces, and finite-sum closure combines the four pieces
\[
\text{positive arm},\qquad
\text{negative arm},\qquad
\text{zero-distance arm},\qquad
\text{outside support}.
\]
On the negative arm the factors \((-1)^k\) are absorbed into the clipped coefficient vector, and the zero-distance and outside-support terms are degree-zero radial residual terms.  A pullback then ties the ball center, radial center, arm, coefficient vector and coordinate to the same index \(\iota\).  These finitely many product, sum, envelope-enlargement and pullback operations preserve polynomial empirical \(L^2\) entropy, changing only the constants; after the final envelope enlargement from \((2\cdot4^p+2)\{B+(p+1)R\}\) to \(C_0B\), the witnesses are the displayed \(A\) and \(v\).
\end{itemize}
The three witnesses \(A\), \(v\), \(C_0\) are chosen before \(\nu\), \(P\), \(h\) and \(B\), which is what makes the final constant uniform in the moment exponent.
% lean: winsorizedScore_hasVCUniformEntropy_all_nu

The parameter-regime clause of \(\mathcal P_{12}(p,\nu,L)\) gives \(L\geq4\), so \(h\leq L^{-1}<1\leq B\) and \(h<B\); the certificate therefore applies at the bandwidth and clipping level of the statement.

\textbf{Step 6: countable reduction.}  Equip \(\mathcal I_{P,h}\) with the product topology inherited from \([h,2h)\subset\mathbb R\), the discrete arm and coordinate sets, the relative topology on \(\mathcal B_P\subset\mathbb R^2\), and the Euclidean topology on \(\mathbb R^{p+1}\).  Choose a fixed countable dense sequence \((\kappa_m)_{m\geq1}\) in this space.  Let
\[
S=\{z:X(z)\in\mathcal X_P,\ X(z)\notin\mathcal B_P\}.
\]
The set \(S\) has \(P\)-probability one: \(X\in\mathcal X_P\) almost surely, and \(P\{X\in\mathcal B_P\}=0\) because the interface has finite one-dimensional Hausdorff measure by \cref{obj:def:cty-a1-a2-class} and hence vanishing planar Lebesgue measure under the displayed density.

Fix \(\iota=(q,t,x,b,j)\in\mathcal I_{P,h}\).  For \(\ell\geq1\) set
\[
\delta_\ell={2h-q\over 16\ell}>0
\]
and consider the target point \((q+4\delta_\ell,t,x,b,j)\in\mathcal I_{P,h}\).  By density of \((\kappa_m)\), choose \(m_\ell\) so that \(\kappa_{m_\ell}=(q_\ell,t_\ell,x_\ell,b_\ell,j_\ell)\) satisfies
\[
\lVert x_\ell-x\rVert_2<\delta_\ell,
\qquad
|q_\ell-(q+4\delta_\ell)|<\delta_\ell,
\qquad
|b_\ell(k)-b(k)|<\delta_\ell\quad(0\leq k\leq p),
\]
with \(t_\ell=t\) and \(j_\ell=j\).  Then \(\kappa_{m_\ell}\to\iota\), and the one-sided construction gives
\[
q+\lVert x_\ell-x\rVert_2<q_\ell
\qquad\text{for every }\ell .
\]
Thus the approximating support bandwidth converges to the target bandwidth while staying above it by enough to cover the center perturbation.  For every \(z\in S\), this slack preserves the closed-kernel decision when \(|D^{\pm}_{P,x}(z)|\leq q\), including the endpoint case \(|D^{\pm}_{P,x}(z)|=q\); when \(|D^{\pm}_{P,x}(z)|>q\), ordinary continuity of distance and of \(q_\ell\) makes the approximating kernel decision eventually agree with the target one.  Since \(z\in\mathcal X_P\setminus\mathcal B_P\), the assignment partition in \cref{obj:def:cty-a1-a2-class} fixes the signed-arm indicator along the approximating boundary centers, and continuity handles the polynomial and clipped-coefficient factors.  Therefore
\[
g_{\kappa_{m_\ell}}(z)\longrightarrow g_\iota(z),
\qquad z\in S .
\]
The class is dominated by the integrable envelope
\[
z\mapsto 2^p\left(|Y(0)|+|Y(1)|+(p+1)2^pR\right),
\]
so dominated convergence also gives convergence of the corresponding population integrals.  Consequently, for every \(n\), almost surely under \(P^n\),
\[
\sup_{\iota\in\mathcal I_{P,h}}
\left|
{1\over n}\sum_{i=1}^n g_\iota(Z_i)-\int g_\iota\,dP
\right|
=
\sup_{m\geq1}
\left|
{1\over n}\sum_{i=1}^n g_{\kappa_m}(Z_i)-\int g_{\kappa_m}\,dP
\right| .
\]
% lean: winsorizedScore_hasCountableEmpiricalSupReduction_at

\textbf{Step 7: the variance-adaptive maximal inequality.}  We use the following countable-reduction form of the variance-adaptive VC maximal inequality.  Let \(\mu\) be a probability law and \(\{G_\iota:\iota\in\mathcal I\}\) a real-valued class.  Suppose that \(0<\sigma<U\), \(A\geq\exp(1)\), \(v\geq1\), every \(G_\iota\) is measurable, \(|G_\iota|\leq U\), \((\int G_\iota^2\,d\mu)^{1/2}\leq\sigma\), every countable subfamily has empirical \(L^2\) covers of size at most \((A/\varepsilon)^v\) at radius \(\varepsilon U\) for \(0<\varepsilon\leq1\), and a countable subfamily realizes the centered empirical supremum almost surely under every finite product law.  Then, with
\[
\ell_{\mu}=\log\!\left(\max\left\{\exp(1),{AU\over\sigma}\right\}\right),
\]
one has, for every \(n\geq1\),
\[
\mathbb E_{\mu^n}^{*}\!\left[
\sup_{\iota\in\mathcal I}
\left|{1\over n}\sum_{i=1}^n G_\iota(W_i)-\int G_\iota\,d\mu\right|
\right]
\leq
c_{\mathrm{VC}}\left\{
\sigma\sqrt{{v\ell_{\mu}\over n}}
+{vU\ell_{\mu}\over n}
\right\},
\qquad c_{\mathrm{VC}}=16384 .
\]
Indeed, choose a reducing sequence \((G_m)_{m\geq1}\).  The envelope makes the countable supremum measurable, integrable, and bounded by \(2U\), so the outer expectation of the continuum supremum is bounded by the ordinary expectation of the countable supremum.  Countable-class symmetrization bounds that expectation by twice the Rademacher complexity.  The Rademacher bound applies Dudley chaining to the empirical \(L^2\) polynomial covers, uses the population \(L^2\) radius \(\sigma\), and controls the empirical radius by symmetrizing the squared bounded class; the resulting quadratic inequality yields the displayed leading \(\sigma\)-term and second-order \(U\)-term, with the numerical constant rounded to \(16384\).

Apply this inequality with \(\mu=P\), \(G_\iota=g_\iota\), \(U=C_0B\), \(\sigma=C_0h\), and with the entropy and countability witnesses from Steps 5 and 6.  Writing
\[
\ell=
\log\left(\max\left\{\exp(1),\ {A\,C_0B\over C_0h}\right\}\right),
\]
we obtain, for every \(n\geq1\),
\[
\mathbb E_{P^n}^{*}\!\left[
\sup_{\iota\in\mathcal I_{P,h}}
\left|
{1\over n}\sum_{i=1}^n g_\iota(Z_i)-\int g_\iota\,dP
\right|
\right]
\leq
c_{\mathrm{VC}}
\left\{
C_0h\sqrt{{v\,\ell\over n}}
+
{v\,C_0B\,\ell\over n}
\right\},
\qquad
c_{\mathrm{VC}}=16384 .
\]
% lean: vcExpectedMaximalInequality_explicit

\textbf{Step 8: calibrating the logarithm.}  Since \(L\geq4\), \(h\leq L^{-1}\) and \(B\geq1\),
\[
4h\leq Lh\leq1\leq B,
\qquad\text{so}\qquad
{B\over h}\geq4
\qquad\text{and}\qquad
\log(B/h)\geq1 .
\]
Also \(A\geq\exp(1)\) gives \(\log A\geq1\), so \(A(B/h)\geq\exp(1)\) and the maximum in \(\ell\) is attained at the second entry.  Putting \(k=1+\log A\),
\[
\ell=\log A+\log(B/h)\leq k\log(B/h),
\]
because \(\{\log A-1\}\{\log(B/h)-1\}\geq0\).
% lean: cty_winsorized_score_maximal_bound

\textbf{Step 9: conclusion.}  Abbreviate
\[
\Gamma_1=\sqrt{{\log(B/h)\over nh^{2}}},
\qquad
\Gamma_2={B\log(B/h)\over nh^{2}},
\]
both nonnegative.  Multiplying the two terms of Step 7 by \(h^{-2}\) and using Step 8,
\[
{1\over h^{2}}\,C_0h\sqrt{{v\,\ell\over n}}
\leq
C_0\sqrt{vk}\;{1\over h}\sqrt{{\log(B/h)\over n}}
=
C_0\sqrt{vk}\,\Gamma_1,
\qquad
{1\over h^{2}}\cdot{v\,C_0B\,\ell\over n}
\leq
C_0vk\,\Gamma_2 .
\]
Therefore, chaining Steps 4, 6, 7 and the two displays above and setting
\[
C=1+\left|c_{\mathrm{VC}}C_0\left\{\sqrt{vk}+vk\right\}\right|,
\]
we obtain
\[
\mathbb E_{P^n}^{*}\!\left[
\sup_{t\in\{0,1\},\,x\in\mathcal B_P}
\left\lVert
\operatorname{score}^{B}_{t,x,h}
-\mathbb E_{P^n}\!\left[\operatorname{score}^{B}_{t,x,h}\right]
\right\rVert
\right]
\leq
c_{\mathrm{VC}}C_0\left\{\sqrt{vk}+vk\right\}(\Gamma_1+\Gamma_2)
\leq
C\left\{
\sqrt{\frac{\log(B/h)}{n h^2}}
+
\frac{B\log(B/h)}{n h^2}
\right\}.
\]
The constant \(C\) is positive and, since \(R\), \(C_0\), \(A\) and \(v\) were fixed from \(p\) and \(L\) alone, depends only on \(p\) and \(L\).
% lean: cty_winsorized_score_maximal_bound
\end{proof}
